%% kapa: exam=1 round=33 year=2022 period=2 form=A subject=accounting subject_ko=회계학 seq=0 kind=subject-instructions
%% source: Q-Net 감정평가사 기출문제 (cst00309) artlSeq=5209997 file=raw/1008/5209997/2022년도 제33회 감정평가사 1차시험 2교시(감정평가 관계법규, 회계학).pdf
%% 과목 공통 지시문(문항 앞에 제시된 안내). 묶음 PDF 에서 첫 문항 앞에 출력된다.
\begin{framedtext}
※ 아래의 문제들에서 특별한 언급이 없는 한 기업의 보고기간(회계기간)은 매년 1월 1일부터 12월 31일까지이다. 또한, 기업은 주권상장법인으로 계속해서 한국채택국제회계기준(K-IFRS)을 적용해오고 있다고 가정하고, 답지항 중에서 물음에 가장 합당한 답을 고르시오. 단, 자료에서 제시한 모든 항목과 금액은 중요하며, 자료에서 제시한 것 이외의 사항은 고려하지 않고 답한다. 예를 들어, 법인세에 대한 언급이 없으면 법인세 효과는 고려하지 않는다. \par
\end{framedtext}
